\chapter{Literature Review}
\label{ch:2}
This chapter first sets out the concepts the rest of the thesis relies on, then reviews four lines of work, and ends with the gap they leave.

\section{Preliminaries}
\label{sec:prelim}

\subsection{Knowledge distillation}
A teacher and a student produce logits $z_t(x), z_s(x) \in \mathbb{R}^C$ for a text $x$ and $C$ classes.
Softened distributions $p^{(T)}_t = \mathrm{softmax}(z_t/T)$ and $p^{(T)}_s = \mathrm{softmax}(z_s/T)$
are compared at a temperature $T > 1$, and the student minimises
\begin{equation}
\mathcal{L}_{\mathrm{KD}} = \beta\, \mathrm{CE}\big(\mathrm{softmax}(z_s), y\big)
 + \alpha\, T^2\, \mathrm{KL}\big(p^{(T)}_t \,\|\, p^{(T)}_s\big)
\end{equation}
\cite{hinton2015distilling}. The temperature spreads probability over the classes the teacher considers
wrong, and the factor $T^2$ keeps the gradient of the second term comparable across temperatures. These
\emph{soft labels} carry how the teacher ranks the wrong classes, which a one-hot label cannot: a teacher
that gives a tweet 0.70 for one class and 0.20 for another tells the student that the two are confusable
there. The argument of this thesis turns on when that information exists. If a teacher has fitted a text
so well that it assigns the gold class a probability near one, its soft label and the gold label are the
same thing.

\subsection{Multi-teacher distillation}
With $K$ teachers the target is a mixture $\bar p = \sum_k w_k\, p^{(T)}_k$. The weights can be uniform,
fixed per teacher, or computed per instance from each teacher's confidence or from its loss against the
gold label, so that the student trusts each teacher where it is reliable \cite{you2017learning,
wu2021one, zhang2022confidence}. Intermediate representations can be matched as well, through learned
projections from the student's width to each teacher's \cite{romero2015fitnets, wu2021one}.

\subsection{The transfer set}
The text on which the student imitates the teacher is called the \emph{transfer set}
\cite{hinton2015distilling}. When it is the labelled split the teachers were fine-tuned on, we call the
distillation \emph{in sample}; when it contains text outside every teacher's training split,
\emph{out of sample}. Section~\ref{sec:3.5} defines both precisely.

\subsection{Measures used throughout}
\emph{Macro-F1} is the unweighted mean of the per-class F1 scores, so that small classes count as much
as large ones. \emph{ROC-AUC} is the probability that a randomly chosen positive is ranked above a
randomly chosen negative; 0.5 is chance and 1.0 is perfect, and it needs no decision threshold.
\emph{Cohen's kappa} measures agreement between two classifiers beyond what chance would produce.
\emph{Expected calibration error} measures how far predicted probabilities are from observed accuracy.
A \emph{paired bootstrap} resamples the test set with replacement many times, recomputes the difference
between two systems on each resample, and reports the central 95\% of those differences as an interval;
a difference whose interval excludes zero is unlikely to be an accident of the particular test set
\cite{koehn2004statistical, dror2018hitchhiker}.

\section{Cyberbullying, Toxicity, and Abuse Carried by Implication}
\label{sec:2.1}
Surveys frame cyberbullying detection as text classification with three persistent problems: scarce
and inconsistent data, weak evaluation practice, and poor transfer across platforms
\cite{rosa2019automatic, salawu2020approaches, emmery2021current}. The wider abusive-language
literature documents the same issues, including training-data quality and its effect on
generalisation \cite{fortuna2018survey, vidgen2020directions, yin2021towards}. The benchmarks used
here are the fine-grained cyberbullying tweet corpus \cite{wang2020sosnet} and, for implicit and
ironic abuse, the Implicit Hate Corpus \cite{elsherief2021latent} and ISHate
\cite{ocampo2023indepth}; ToxiGen generates implicit examples adversarially
\cite{hartvigsen2022toxigen} and HateXplain adds rationales \cite{mathew2021hatexplain}. Sarcasm
resources are iSarcasm and iSarcasmEval \cite{oprea2020isarcasm, abufarha2022semeval} and the
SemEval-2018 irony task \cite{vanhee2018semeval}. Domain-specialised encoders include HateBERT
\cite{caselli2021hatebert}, the TweetEval models \cite{barbieri2020tweeteval} and the ToxDect
RoBERTa \cite{zhou2021challenges}. Earlier datasets and the pitfalls of duplicate-heavy corpora are
documented in \cite{vanhee2018automatic, founta2018large, davidson2017automated, borkan2019nuanced}.
The unlabelled transfer set of Section~\ref{sec:4.4} takes its text from Davidson et al.
\cite{davidson2017automated}, OLID \cite{zampieri2019olid} and HatEval \cite{basile2019hateval}, and
excludes Founta et al. under the provenance rule of Section~\ref{sec:4.5}.

A separate line of work argues that the hard case is not offensive vocabulary but its absence.
ElSherief et al. \cite{elsherief2021latent} build a taxonomy of implicit hate and a corpus in which
implication is annotated as such, and report that models competent on explicit hate lose most of
their accuracy on it; Ocampo et al. \cite{ocampo2023indepth} separate hate speech into explicit and
implicit layers and show the implicit layer is where detectors and annotators both struggle. The
annotation literature adds that the texts annotators disagree on are disproportionately the ones
whose hostility is implied \cite{uma2021learning, davani2022dealing, plank2022problem,
peterson2019human}. Two things follow for a thesis about compact models. A corpus that does not
label implication cannot be used to train for it, and the domain-specialised encoders a distillation
committee would recruit were themselves trained on corpora of that kind. We are not aware of work
that asks whether the ability to read implication can be moved into a deployable model, or by what
mechanism; Section~\ref{sec:5.9} asks it and finds no intervention in this grid that moves it, and
Section~\ref{sec:4.2} measures how far the label itself can be trusted.
\section{Distilling Compact Language Models, and the Transfer Set}
\label{sec:2.2}
Language models have grown far past what a moderation pipeline can run on every post
\cite{brown2020language}, which is the pressure distillation answers. It transfers a large model's
softened outputs to a smaller one \cite{hinton2015distilling, bucilua2006model};
intermediate-feature transfer follows FitNets \cite{romero2015fitnets}; surveys organise response-,
feature- and relation-based variants \cite{gou2021knowledge}. For BERT-style encoders the line runs
through DistilBERT \cite{sanh2019distilbert}, Patient KD \cite{sun2019patient}, TinyBERT
\cite{jiao2020tinybert}, MiniLM \cite{wang2020minilm} and MobileBERT \cite{sun2020mobilebert}.

Hinton et al. named the data the student imitates the teacher on the \emph{transfer set}, and
observed that it need not be the labelled training data \cite{hinton2015distilling}. The recipes
that report the largest gains for compact students all use one: Tang et al. distil BERT into a
BiLSTM on augmented text \cite{tang2019distilling}, TinyBERT augments the task data
\cite{jiao2020tinybert}, and Turc et al. show that pre-training the compact student matters more
than the distillation recipe and that the size and properties of the unlabelled transfer data are
variables in their own right \cite{turc2019wellread}. Distillation does not always work as expected
\cite{stanton2021does, cho2019efficacy}: students fail to match teachers they have the capacity to
match, and the fidelity of the imitation depends on the data it is practised on. Beyer et al. make
the same point from the other side, casting distillation as function matching that needs the
teacher's function sampled on many inputs, consistently, over long schedules
\cite{beyer2022knowledge}. Our results are a domain-specific instance of these observations,
measured with the controls they call for: the transfer set's size, its composition, the number of
updates, the labeller, and the form of the label are each varied with the rest held fixed.
\section{Multi-Teacher and Adaptive Distillation}
\label{sec:2.3}
Ensembles of teachers \cite{you2017learning, fukuda2017efficient} are combined with fixed or learned
weights: adaptive multi-level weighting \cite{liu2020adaptive}, gradient-space adaptation
\cite{du2020agree}, reinforced teacher selection \cite{yuan2021reinforced} and confidence-aware
weighting \cite{zhang2022confidence}. For language models, MT-BERT co-fine-tunes several teachers
and weights their soft labels by prediction error while aligning hidden states through learned
projections \cite{wu2021one}; dynamic weighting by teacher confidence has been used for semantic
parsing \cite{zou2025dynamic}. Passalis et al. model information flow explicitly when teacher and
student architectures differ \cite{passalis2020heterogeneous}; the heterogeneous student here is
trained without such a model (Section~\ref{sec:3.3}).

The method we audit belongs to the error-weighted family: its per-instance weights follow
\cite{wu2021one, zhang2022confidence}, and the combination of error-weighted soft labels with
projected hidden-state alignment is MT-BERT's. We do not claim that combination as novel. What the
family shares, and what none of the papers reviewed here tests, is a pair of assumptions: that the
teachers collectively know more than any one of them in a way their outputs reveal, and that
reliability scored against the gold label on the training data still means something once the
teachers have been fine-tuned on that data. In the abusive-language literature the claim that
several teachers beat one is typically reported without a no-distillation control, without the
teachers' own scores, and without seeds or intervals \cite{prasomphan2025mtkd}.
Section~\ref{sec:5.5} tests both assumptions and finds both false on a committee assembled in the
usual way, and Section~\ref{sec:5.7} finds the committee no better a labeller than one teacher out
of sample either.
\section{Evaluation Practice}
\label{sec:2.4}
Targeted functional tests expose what held-out accuracy hides \cite{rottger2021hatecheck}; our probe
sets and threshold-free AUC (Section~\ref{sec:4.6}) are in that spirit, built by rule from published
corpora and awaiting the human validation HateCheck's authors gave theirs. Reporting practice
matters as much as the tests: papers should state seeds, search budget and compute
\cite{dodge2019show}, test differences with paired bootstrap \cite{koehn2004statistical,
dror2018hitchhiker}, and know the minimum effect their test set can detect, which for most NLP
benchmarks is larger than the differences reported \cite{card2020power}. On our 4,326-post test set
a paired interval is about 0.014 wide, so a difference under 0.007 macro-F1 is undetectable; every
claim in this thesis is read against that.

\section{Summary and Research Gap}
\label{sec:gap}
Table~\ref{tab:gap} summarises what each line of work establishes, what it leaves open for this thesis,
and where the thesis addresses it.

\begin{table}[htbp]
\centering
\small
\caption{What the literature establishes and what it leaves open.}
\label{tab:gap}
\begin{tabular}{p{3.2cm}p{4.3cm}p{4.3cm}p{2.4cm}}
\toprule
Line of work & Establishes & Leaves open & Addressed in \\
\midrule
Cyberbullying and toxicity detection & Benchmarks and strong fine-tuned encoders & Compact models are
evaluated on aggregate F1 only; implicit abuse is not labelled & Chapter~\ref{ch:4}, Section~\ref{sec:5.9} \\
Implicit hate & Implicit hate needs its own label and benchmark & Whether a compact model can learn it by
distillation & Sections~\ref{sec:5.5} and~\ref{sec:5.9} \\
Distilling compact models & Transfer sets and pre-training matter & Not measured in this domain with the
confounds controlled & Sections~\ref{sec:5.7} and~\ref{sec:5.8} \\
Multi-teacher weighting & Several teachers are claimed to beat one & The no-teacher control, seeds and
intervals; reliability is read in sample & Sections~\ref{sec:5.2} to~\ref{sec:5.6} \\
Evaluation practice & Functional tests, reporting standards, power & A threshold-free measure of
implication for compact detectors & Sections~\ref{sec:4.6} and~\ref{sec:5.9} \\
\bottomrule
\end{tabular}
\end{table}
